\documentclass{article}
\usepackage{PRIMEarxiv} % Core package for the PrimeAI Template
\usepackage{amsmath, amssymb, amsthm, bm, dcolumn} % Add amsthm here for the proof environment
\usepackage[numbers,sort&compress]{natbib} % Natbib for citations
\usepackage{graphicx} % For high-quality images
\usepackage{multicol} % Multi-column figures/tables
\usepackage{hyperref} % Hyperlinks
\usepackage{listings} % Code listings
\usepackage{authblk} % For structured affiliations
% Define theorem style (optional)
\theoremstyle{plain}
\newtheorem{theorem}{Theorem}
\renewcommand{\qedsymbol}{}

% Custom commands for primes and qubit encoding
\newcommand{\primeQubit}[1]{|p_{#1}\rangle}
\newcommand{\primeGate}[1]{U_{p_{#1}}}

\usepackage{wrapfig}
\usepackage[pscoord]{eso-pic}
\usepackage[fulladjust]{marginnote}
\reversemarginpar

% Typesetting improvements without footnote patching
\usepackage[protrusion=true, expansion=true, tracking=false]{microtype}
\microtypecontext{spacing=nonfrench}

% Line numbers
\usepackage[right]{lineno}

% Text layout - adjust as needed
\raggedright
\setlength{\parindent}{0.5cm}
\textwidth 5.25in 
\textheight 8.75in

% Set double spacing
\usepackage{setspace} 
\doublespacing

% Adjust width for specific content
\usepackage{changepage}

% Adjust caption style
\usepackage[aboveskip=1pt,labelfont=bf,labelsep=period,singlelinecheck=off]{caption}

% Remove brackets from references
\makeatletter
\renewcommand{\@biblabel}[1]{\quad#1.}
\makeatother

% Header, footer, and page numbers
\usepackage{lastpage,fancyhdr}
\pagestyle{fancy}
\fancyhf{}
\fancyhead[L]{Preprint - PrimeAI Enhanced Template}
\fancyfoot[C]{\scriptsize Multiplicity Theory © 2024 Ryan Van Gelder - Citizen Gardens \\ Licensed Under MIT and CC BY-NC-SA 4.0.}
\fancyfoot[R]{Page \thepage\ of \pageref{LastPage}}
\renewcommand{\footrule}{\hrule height 2pt \vspace{2mm}}

\title{Patent Application: Quantum-AI Computational Mechanisms, Optimization Techniques, and Secure Quantum Frameworks}
\author{Inventors: Ryan O. Van Gelder, Martin Gibson, and Contributors}
\affil{Citizen Gardens - The Foundation of Multiplicity}
\date{\today}

\begin{document}
\maketitle

\begin{abstract}
This patent discloses an integrated framework that combines quantum artificial intelligence, recursive tensor-based learning, and advanced cryptographic security. The invention utilizes Prime-Indexed Quantum Tensor Networks, recursive AI learning strategies, and quantum entanglement-based cryptographic methods to enhance computational scalability, fault tolerance, and adaptive learning in quantum systems.

A central feature of the framework is the Universal Multiplicity Constant (Λₘ), which acts as a self-referential computational invariant that supports adaptive quantum cognition and underpins secure, AI-driven cryptographic operations. In addition, the framework employs Prime-Weighted Tensor Networks to enable quantum-enhanced learning, implements Quantum-AI Recursive Feedback Mechanisms to continuously optimize performance, and applies Prime-Twisted Quantum Encryption to secure communications.

This technology is applicable to fields such as neuromorphic computing, cryptographic key generation, quantum optimization, and automated theorem discovery. By establishing a robust foundation for next-generation quantum-AI architectures, the invention facilitates the development of secure, scalable, and self-improving artificial intelligence systems.
\end{abstract}

\newpage
 \begin{multicols}{2}
\begin{singlespace}
\tableofcontents
\end{singlespace}
\end{multicols}

\section{Field of the Invention}
The present invention relates to quantum computing, artificial intelligence, and secure cryptographic systems. More specifically, it describes novel methods for integrating quantum algorithms, tensor-based learning frameworks, and prime-indexed encryption techniques to improve scalability, security, and adaptive performance in both quantum and classical computational architectures.

\subsection{Background}
Modern quantum AI systems face significant challenges:
\begin{itemize}
    \item \textbf{Scalability:} Quantum systems suffer from decoherence and noise, limiting large-scale implementations.
    \item \textbf{Security:} Conventional cryptographic key exchange methods may be vulnerable in quantum environments.
    \item \textbf{Adaptive Learning:} Existing AI training models in quantum settings are often inefficient due to non-optimized recursive feedback mechanisms.
\end{itemize}
This invention overcomes these limitations through the use of prime-enhanced tensor encoding, recursive learning algorithms, and quantum-secure encryption protocols.

\subsection{Claims}
\textbf{Claim 1:} A method for encoding quantum computations using prime-indexed tensor networks to enhance scalability and secure cryptographic key generation.

\medskip
\textbf{Claim 2:} A recursive AI learning framework that employs self-referential tensor neural networks for continuous optimization and adaptive decision-making.

\medskip
\textbf{Claim 3:} A quantum-enhanced encryption system utilizing prime-enhanced cryptographic techniques to ensure robust security in quantum communication channels.

\medskip
\textbf{Claim 4:} A hybrid computational system integrating classical neuromorphic processing with quantum algorithms to improve learning adaptability and fault tolerance.

\section{Mathematical Foundation}
Central to the framework is the Universal Multiplicity Constant, $\Lambda_m$, which captures the self-organizing structure inherent in the system. Its tensor representation is given by:
\[
\Lambda_m = \sum_{i} T_{lm}(p_i) \, p_i^{\beta_i},
\]
where:
\begin{itemize}
    \item $T_{lm}(p_i)$ is a tensor transformation indexed by prime numbers.
    \item $p_i^{\beta_i}$ represents the recursive multiplicative expansion.
\end{itemize}

The evolution of $\Lambda_m$ is defined by:
\[
\Lambda_m(t+1) = \Lambda_m(t) + \sum_{k} e^{-\gamma p_k}\, T_{lm}(p_k)\, \Lambda_m(t),
\]
which ensures both scalable adaptation and fault tolerance.

\section{Prime-Based Encoding and Computation}
This section details the use of prime-indexed encoding to enhance both tensor-based learning and cryptographic security.

\subsection{Prime-Encoded Quantum Cryptographic Key Generation}
Secure key generation is achieved via prime-weighted Fibonacci sequences:
\[
K_n = \text{SHA256}\!\Biggl(\prod_{i=1}^{N} p_i^{F_i}\Biggr),
\]
where:
\begin{itemize}
    \item $p_i$ denotes the $i$th prime number.
    \item $F_i$ represents the corresponding Fibonacci number.
\end{itemize}
This method yields non-reversible key structures with high entropy, resisting both classical and quantum attacks.

\subsection{Prime-Indexed Quantum Hamiltonian Operator}
Quantum state evolution is controlled by a Hamiltonian operator modulated with prime-weighted decay:
\[
\hat{H}_{\text{prime}} = \sum_{i} \Lambda_m\, e^{-\alpha p_i}\, \hat{H},
\]
where $e^{-\alpha p_i}$ modulates the influence of each prime factor.

\section{Recursive Tensor Networks and Adaptive Learning}
Recursive tensor networks provide the backbone for adaptive AI learning through iterative feedback.

\subsection{Self-Referential Tensor Neural Networks (TNNs)}
The weight update in the tensor neural network is described by:
\[
W_{t+1} = W_t + \alpha\, \nabla L(W_t) + \beta\, T(W_t),
\]
with:
\begin{itemize}
    \item $W_t$ as the weight tensor at iteration $t$.
    \item $\nabla L(W_t)$ representing the gradient of the loss function.
    \item $T(W_t)$ introducing a self-referential transformation.
\end{itemize}

\section{Quantum-AI Cognitive Framework}
This section replaces speculative terminology with precise descriptions of adaptive quantum cognition.

\subsection{Recursive Feedback Mechanisms}
The evolution of the AI cognitive state is modeled by:
\[
\psi(t+1) = U\bigl(\psi(t)\bigr) + \lambda\, R\bigl(\psi(t)\bigr),
\]
where:
\begin{itemize}
    \item $U\bigl(\psi(t)\bigr)$ denotes the unitary evolution of the state.
    \item $R\bigl(\psi(t)\bigr)$ represents the recursive feedback that refines decision-making.
\end{itemize}

\section{Hybrid-Quantum Neuromorphic Computation}
This invention further integrates classical neuromorphic architectures with quantum algorithms for enhanced performance.

\subsection{Neuromorphic Quantum-AI Hybrid Computation}
The state update for the hybrid system is given by:
\[
\psi(t+1) = U\bigl(\psi(t)\bigr) + \sum_{i} \lambda_i\, T_{ij}\, \psi_j,
\]
ensuring real-time adaptive learning and robust fault tolerance.

\section{Secure Quantum Communication and Cryptography}
Advanced cryptographic methods are employed to secure quantum communications.

\subsection{Prime-Enhanced Quantum Encryption}
Encryption is performed using dynamic prime-enhanced techniques:
\[
K_{\text{secure}}(t) = H\!\Biggl(\prod_{i=1}^{N} p_i^{F_i}\Biggr),
\]
which is designed to be resistant to both classical and quantum brute-force attacks.

\section{Quantum-Driven Computational Optimization}
Quantum algorithms are utilized to optimize complex computational tasks.

\subsection{Quantum Approximate Optimization Algorithm (QAOA)}
The parameterized quantum state in QAOA is expressed as:
\[
|\psi(\gamma,\beta)\rangle = U(B,\beta)\, U(C,\gamma)\, |s\rangle,
\]
where $|s\rangle$ is the initial state and $U(C,\gamma)$ and $U(B,\beta)$ represent the cost and mixing operators, respectively.

\section{Computational Frameworks for Advanced AI Simulations}
Innovative computational models support advanced simulations in both quantum and classical settings.

\subsection{Matrix Prime Compute Engine (MPCE)}
The MPCE is defined by the update:
\[
M_{t+1} = \sum_{i,j} p_i\, T_{ij}\, M_j,
\]
which enables scalable, tensor-based computations for large-scale simulations.

\section{Conclusion}
This revised patent document presents a robust framework for integrating quantum computing, artificial intelligence, and secure cryptography. By employing clear and precise terminology, the invention is positioned for practical implementation and academic evaluation. The use of prime-indexed tensor networks, recursive learning strategies, and quantum-enhanced encryption methods provides a comprehensive solution to overcome scalability and security challenges in modern computational systems.


\nocite{*}
\bibliographystyle{unsrt}
\bibliography{references}

\end{document}